\documentclass[12pt]{article}
\usepackage[utf8]{inputenc}
\usepackage[T1]{fontenc}
\usepackage{hyperref}
\title{Iduna Language Resources for K-12 Curriculum Development}
\author{Compiled by Claude Sonnet 5 (claude-sonnet-5) for Andr\'as Kornai}
\date{}
\begin{document}
\maketitle

\section{Introduction}
Iduna (ISO 639-3 viv; also called Vivigani) is an Austronesian language of Goodenough Island, Milne Bay Province, Papua New Guinea, with about 6,000 speakers (a 1984 figure). Documentation is mostly linguistic and was produced by SIL in the 1960s and 1970s. No school curriculum was found in the search results.

\section{Primary Educational Resources}
A draft ``Vivigani language lessons'' PDF is mentioned on the Wikipedia article. It was not opened, and its link was not seen in search results, so it is not cited here.

\section{Linguistic and Reference Materials}
The following works are listed as references on Wikipedia pages about Iduna. Only those pages were seen.
\begin{itemize}
\item Wikipedia, Iduna language: \url{https://en.wikipedia.org/wiki/Iduna_language}
\item Wikipedia mirror with the ISO 639 entry: \url{https://wikipedia.com/wiki/ISO_639:viv}
\item Joyce Huckett, ``Notes on Iduna Grammar'' (1974) and ``Iduna Sentence Structure'' (1976), both mentioned in the pages above. No direct links were found.
\item Numerals in Iduna, Max Planck Institute: \url{https://lingweb.eva.mpg.de/channumerals/Iduna.htm}
\item Glosbe Iduna--English pages (community-built and not a scholarly dictionary): \url{https://glosbe.com/viv/en}
\end{itemize}

\section{Religious and Community Resources}
\begin{itemize}
\item Joshua Project, Vivigani/Iduna (its location differs from the linguistic sources and should be checked): \url{https://joshuaproject.net/people_groups/15801}
\end{itemize}

\section{Digital and Audio Resources}
\begin{itemize}
\item Global Recordings Network, Iduna: Goiala language: \url{https://globalrecordings.net/en/language/10791}
\item Global Recordings Network, Iduna: Ufaufa language: \url{https://globalrecordings.net/en/language/10793}
\end{itemize}
These two pages are labelled as Iduna varieties. The search did not confirm what they contain.

\section{Curriculum Development Context}
The language is spoken by all in the community, which is a good basis for mother-tongue schooling. However, no primers, textbooks or school programmes were found. A curriculum would need new material, built on the SIL grammar and phonology works and developed with local speakers.
\end{document}